% Source pp. 214-226 (Lectures19-25.pdf).
\section{Wilson-Fisher Fixed Point}

Previously, we have studied the simplest example of fixed point, the Gaussian
fixed point corresponding to free massless scalar field. It is easy to argue,
however, that the Gaussian fixed point cannot be responsible for the ordinary
transition in a magnet in three dimensions (the Gaussian fixed point describes the
so called ``tricritical behavior'', as we will briefly discuss later). Indeed, it
was very important for the arguments given in the previous lecture that the
critical surface of $\mathcal A_*$ has codimension one in $\Sigma_{\text{symm}}$,
i.e.\ that there is only one relevant field which respects the symmetry of the
problem ($\varphi\to-\varphi$ in our case). However, as we have seen, the gaussian
fixed point has more than one symmetric relevant field for all $d<4$. In
particular, for $d=3$ two fields, $\Phi_2\sim\varphi^2$ and $\Phi_4\sim\varphi^4$,
are relevant. Therefore there must be another fixed point describing the
universality class of the magnetic phase transition.

The best theoretical evidence for the existence of another nontrivial fixed point
at $d=3$ (on a formal level, at all $2<d<4$) was given by Wilson and Fisher, and
it is usually referred to as the Wilson-Fisher fixed point $\mathcal A_{*WF}$.

At $d<4$ the topology of the RG flow relating the fixed points $\mathcal A_{*G}$
and $\mathcal A_{*WF}$ is likely to be the following. The unstable manifold
$U(\mathcal A_{*G})$ at $d<4$ has dimensionality two; it is generated by the two
relevant fields $\Phi_2$ and $\Phi_4$ (again, we restrict attention to the
symmetric space $\Sigma_{\text{symm}}$). The fixed point $\mathcal A_{*WF}$ is
likely to lie in $U(\mathcal A_{*G})$, so that the RG flow on this unstable
manifold would look qualitatively as shown in figure~\ref{fig:wfflow}.
\begin{figure}[ht]
\centering\figfile{wfflow}
\caption{RG flow on the unstable manifold $U(\mathcal A_{*G})$ at $d<4$. The
trajectory $C$ along the critical surface $\Sigma_c$ ($M^2=0$) connects the
Gaussian fixed point to the Wilson-Fisher fixed point.}
\label{fig:wfflow}
\end{figure}

According to this picture there is a trajectory $C$ connecting the fixed point
$\mathcal A_{*G}$ to $\mathcal A_{*WF}$. Therefore $\mathcal A_{*G}$ lies on the
critical surface $\Sigma_c$ which also passes through $\mathcal A_{*WF}$. The
unstable manifold $U(\mathcal A_{*WF})$ is then one-dimensional, as expected.

Note that the whole unstable manifold $U(\mathcal A_{*G})$ in this picture, being
an element of $\Sigma(\infty)$, represents local field theories. Namely, it
represents two-parameter family of $\varphi^4$ field theories. These field
theories can be characterized by two parameters, the renormalized mass parameter
$M^2$ and the renormalized coupling constant $\lambda$. Therefore these parameters
can be taken for coordinates in the above picture. Although there is an ambiguity
in the choice of these coordinates (the scheme dependence), the critical surface
corresponding to $M^2=0$ is defined by the scheme-independent condition
\begin{equation}\label{eq:masslesscond}
\Gamma^{(2)}(p^2)\big|_{p^2=0}=0\,.
\end{equation}
In these coordinates the fixed point $\mathcal A_{*WF}$ lies at $M^2=0$ and
$\lambda=\lambda_*$, where the value of $\lambda_*$ depends on the scheme.

We can formally consider the $\varphi^4$ theory in $d$ dimensional space, treating
$d$ as continuous parameter. What happens when $d$ approaches 4? We know that at
$d=4$ the field $\varphi^4$ becomes marginally irrelevant, and for $d>4$ it is
strictly irrelevant, and so for all $d\ge4$ the gaussian fixed point has only one
relevant field $\varphi^2$. This behavior emerges naturally if one assumes that
when $d$ approaches the value 4 from below the fixed point $\mathcal A_{*WF}$ gets
very close to the gaussian fixed point, and at $d=4$ these two fixed points merge.
At $d>4$ $\lambda_*$ becomes negative, and it does not correspond to field theory
in usual sense.\footnote{The ``magnetic criticality'' at $d\ge4$ then is described
by the gaussian fixed point.} This idea leads to the patterns of the RG flow at
$d=4$ and at $d>4$ shown in figure~\ref{fig:wfmerge}.
\begin{figure}[ht]
\centering\figfile{wfmerge}
\caption{The RG flow near the Gaussian fixed point at $d=4$ (left), where
$\varphi^4$ is marginally irrelevant, and at $d>4$ (right), where the
Wilson-Fisher fixed point has moved to negative $\lambda$.}
\label{fig:wfmerge}
\end{figure}

Note that the pattern at $d=4$ is exactly what one expects to see when $\varphi^4$
is marginally irrelevant, in which case we know that the flow in the vicinity of
the gaussian fixed point is described by the equations
\begin{align*}
\frac{\dd}{\dd l}\,\lambda&=-\beta(\lambda)=-3\,\frac{\lambda^2}{(4\pi)^2}+O(\lambda^3)\,,\\
\frac{\dd}{\dd l}\,M^2&=\bigl(2-\gamma_{\varphi^2}(\lambda)\bigr)\,M^2=\left(2-\frac\lambda{(4\pi)^2}+O(\lambda^2)\right)\,M^2\,.
\end{align*}
On the basis of these arguments Wilson and Fisher have assumed that if one
formally takes $d=4-\epsilon$ with $\epsilon$ small, say $d=3.99$, the fixed point
$\mathcal A_{*WF}$ can be observed within the framework of renormalized
perturbation theory in $\varphi^4$ field theory.

Let us briefly discuss what changes in the renormalized perturbation theory of
$\varphi^4$ if one takes $d=4-\epsilon$ instead of $d=4$. As before, we start with
the action in its ``bare'' form
\[
\mathcal A=\int \dd^dx\,\left[\frac12\,(\partial\varphi_0)^2+\frac{m_0^2}2\,\varphi_0^2+\frac{\lambda_0}{4!}\,\varphi_0^4\right]
\]
Note that in $d$ dimensions the canonical dimension of $\varphi_0$ is $(d-2)/2$
and hence coupling constant $\lambda_0$ has positive mass dimension,
\[
[\varphi_0]=\frac{d-2}2=1-\frac\epsilon2\,;\qquad[\lambda_0]=4-d=\epsilon\,.
\]
To develop renormalized perturbation theory we rewrite this action in terms of the
renormalized field and renormalized parameters,
\begin{equation}\label{eq:Aren21}
\mathcal A=\int \dd^dx\,\left[\frac12\,(\partial\varphi)^2+\frac{M^2}2\,\varphi^2+\frac\lambda{4!}\,\varphi^4
+\frac{\delta Z}2\,(\partial\varphi)^2+\frac{\delta M^2}2\,\varphi^2+\frac{\delta\lambda}{4!}\,\varphi^4\right]\,,
\end{equation}
where the counterterms are determined from suitably chosen normalization
conditions.

It is worth stressing here that for $d<4$ the diagrams describing the field and
the coupling constant renormalizations do not diverge; the only divergent
renormalization is the one of the mass parameter. Therefore the counterterms
$\delta Z$ and $\delta\lambda$ have finite limit when the cutoff parameter
$\Lambda$ goes to infinity. Therefore at $d<4$ these renormalizations are
introduced not for the purpose of eliminating the divergences, which are absent,
but in order to introduce more suitable coordinate $\lambda$ on
$U(\mathcal A_{*G})$ instead of $\lambda_0$. Although for $d<4$ (unlike the case
of $d=4$) the bare coupling $\lambda_0$ can be used as local coordinate on this
unstable manifold, we will see soon that it is not too useful for the purpose of
studying the fixed point $\mathcal A_{*WF}$.

As the fixed point $\mathcal A_{*WF}$ is expected to be located on the critical
surface, it suffices to study the massless $\varphi^4$ theory which is defined by
the condition \eqref{eq:masslesscond},
\[
\Gamma^{(2)}(p^2)\big|_{p^2=0}=0\,.
\]
As before, choose some normalization scale $\mu$ and assume the following simple
normalization conditions
\begin{equation}\label{eq:norm21}
\frac{\dd}{\dd p^2}\,\Gamma^{(2)}(p^2)\Big|_{p^2=\mu^2}=1\,,\qquad
\Gamma^{(4)}(p_1,\dots,p_4)\Big|_{[(p_i+p_j)^2=4\mu^2]}=\lambda\,.
\end{equation}
The first of them fixes the normalization of $\varphi$ while the second gives the
definition of the renormalized coupling constant $\lambda$.

At $d<4$ the Callan-Symanzik equation for the massless theory is
\begin{equation}\label{eq:CS21}
\left[x\frac\partial{\partial x}+N\,D(\lambda)+\beta(\lambda)\,\frac\partial{\partial\lambda}\right]\bigl\langle\varphi(x_1)\cdots\varphi(x_N)\bigr\rangle_\lambda=0\,;
\end{equation}
its form reflects the general properties of the RG we sketched the last time,
which did not depend on $d$. However for generic $d$
\[
D(\lambda)=\frac{d-2}2+\gamma(\lambda)\,,\qquad\gamma(0)=0\,,
\]
because canonical dimension of $\varphi$ is $(d-2)/2$. Our nearest goal is to
determine the RG functions $\beta(\lambda)$ and $\gamma(\lambda)$ for $d<4$. The
first one is particularly important because the RG flow along the massless part of
$U(\mathcal A_{*G})$ (the trajectory $C$ in figure~\ref{fig:wfflow}) is described
by the equation
\[
\frac{\dd}{\dd l}\,\lambda=-\beta(\lambda)\,.
\]
We can calculate the correlation functions perturbatively, order by order in
$\lambda$, and then find $\beta$ and $\gamma$ from the CS equation \eqref{eq:CS21}
itself, as we did before in the case $d=4$. It is convenient to write the
renormalized coupling constant as
\[
\lambda=\mu^{4-d}\,\bar\lambda\,,
\]
where now $\bar\lambda$ is dimensionless parameter. Then by dimensional counting
the $N$-point correlation functions have the form
\[
\bigl\langle\varphi(x_1)\cdots\varphi(x_N)\bigr\rangle_\lambda=\mu^{N(d-2)/2}\,C^{(N)}(\mu x_1,\dots,\mu x_N|\bar\lambda)
\]
and therefore the CS equation can be rewritten as
\begin{equation}\label{eq:CSbar}
\left[\mu\frac\partial{\partial\mu}+N\,\bar\gamma(\bar\lambda)+\bar\beta(\bar\lambda)\,\frac\partial{\partial\bar\lambda}\right]\bigl\langle\varphi(x_1)\cdots\varphi(x_N)\bigr\rangle_{\bar\lambda}=0\,,
\end{equation}
where $\bar\beta$ and $\bar\gamma$ are the dimensionless functions of the
dimensionless parameter $\bar\lambda$ defined as
\[
\beta(\lambda)=\mu^\epsilon\,\bar\beta(\lambda/\mu^\epsilon)\,,\qquad\gamma(\lambda)=\bar\gamma(\lambda/\mu^\epsilon)\,,
\]
and the subscript $\bar\lambda$ indicates that the correlation function is
expressed through $\bar\lambda$ instead of the dimensional $\lambda$. In what
follows I will omit all these bars and use $\lambda$ for $\bar\lambda$, $\beta$
for $\bar\beta$, etc.; anyhow, these quantities coincide if $\mu$ is taken to be
the unit of mass. The equation \eqref{eq:CSbar} can also be rewritten in terms of
the momentum-space proper vertices,
\begin{equation}\label{eq:CSGamma21}
\left[\mu\frac\partial{\partial\mu}-N\,\gamma(\lambda)+\beta(\lambda)\,\frac\partial{\partial\lambda}\right]\Gamma^{(N)}(p_1,\dots,p_N)=0\,.
\end{equation}

We can now repeat our calculations of the leading perturbative contributions to
the proper vertices, which are close to identical to the calculations within the
dimensional regularization method, and obtain
\[
\Gamma^{(2)}(p^2)=p^2+O(\lambda^2)\,,
\]
and
\begin{multline}\label{eq:G4d}
\Gamma^{(4)}(p_1,\dots,p_4)=\mu^{4-d}\Bigl\{\lambda-\frac12\,\frac{\lambda^2}{(4\pi)^{d/2}}\,\frac{\Gamma(2-d/2)\,\Gamma^2(d/2-1)}{\Gamma(d-2)}\\
\times\left[\left(\frac{p_{12}^2}{\mu^2}\right)^{(d-4)/2}+\left(\frac{p_{13}^2}{\mu^2}\right)^{(d-4)/2}+\left(\frac{p_{14}^2}{\mu^2}\right)^{(d-4)/2}-3\,\bigl(2^{d-4}\bigr)\right]+O(\lambda^3)\Bigr\}\,.
\end{multline}
In the last equation the counterterm contribution is fixed by the normalization
condition \eqref{eq:norm21}, i.e.
\[
\Gamma^{(4)}(p_1,\dots,p_4)\big|_{[(p_i+p_j)^2=4\mu^2]}=\mu^{4-d}\,\lambda\,.
\]
Let me remind you that in dimensional regularization of $d=4$ theory the
counterterm contribution was introduced to cancel the $d=4$ singularity of the
Gamma-function $\Gamma(2-d/2)\sim2/(4-d)$. For $d<4$ the loop diagram contributions
are finite even without any counterterm. We can expand in terms of the bare
coupling constant $\lambda_0$ instead of the renormalized one; in this expansion
the counterterm contribution in \eqref{eq:G4d} would be absent. It is instructive
to remember however that when $\epsilon=4-d$ is small the coefficients of this
expansion in terms of the bare coupling constant are large,
\begin{equation}\label{eq:G4bare}
\Gamma^{(4)}(p_1,\dots,p_4)\sim\lambda_0-\frac1\epsilon\,\frac{\lambda_0^2}{(4\pi)^2}\Bigl((p_{12}^2)^{-\epsilon/2}+(p_{13}^2)^{-\epsilon/2}+(p_{14}^2)^{-\epsilon/2}\Bigr)+\cdots\,.
\end{equation}
At $d=4-\epsilon<4$ the ``coordinate transformation'' from $\lambda_0$ to $\lambda$
is nonsingular in some neighborhood of 0, but at small $\epsilon$ it is ``almost
singular''.\footnote{The radius of convergence of the series $\lambda_0(\lambda)$
is $\sim\epsilon$.}

Substituting \eqref{eq:G4d} into \eqref{eq:CSGamma21} we find
$\gamma(\lambda)=O(\lambda^2)$ and
\[
\beta(\lambda)=(d-4)\,\lambda+3\cdot2^{d-4}\,\frac{\lambda^2}{(4\pi)^{d/2}}\,\frac{\Gamma(3-d/2)\,\Gamma^2(d/2-1)}{\Gamma(d-2)}+O(\lambda^3)\,.
\]
The term $(d-4)\lambda$ appeared because of the factor $\mu^{4-d}$; it represents
the canonical dimension of the coupling constant $\lambda$ in the action
\eqref{eq:Aren21}.

Let us assume now that $d=4-\epsilon$ where $\epsilon$ is small,
\[
\epsilon\ll1\,.
\]
Then
\[
\beta(\lambda)=-\epsilon\,\lambda+3\,\frac{\lambda^2}{(4\pi)^2}+O(\lambda^3)\,.
\]
The linear term $-\epsilon\lambda$ brings in an important difference from the case
$d=4$. For small $\epsilon$ the RG flow equation
\[
\frac{\dd}{\dd l}\,\lambda=-\beta(\lambda)
\]
has nontrivial fixed point
\[
\lambda_*=\frac{(4\pi)^2}3\,\epsilon\,.
\]
The RG flow looks like this:
\[
\diac[0.6]{\draw[line width=0.8pt] (-2.5,0) -- (6.5,0) node[right] {$\lambda$};
  \draw[arr] (-0.8,0) -- (-2.2,0); \draw[arr] (0.8,0) -- (2.4,0); \draw[arr] (5.8,0) -- (4.6,0);
  \fill (0,0) circle (2.4pt) node[above=2pt] {$0$}; \fill (3.6,0) circle (2.4pt) node[above=2pt] {$\lambda_*$};}
\]
This fixed point was obtained in the leading order of the renormalized
perturbation theory. Note however that for small $\epsilon$ the fixed point value
$\lambda_*$ is small. The higher order terms can shift its position by terms
$\sim\epsilon^2$,
\begin{equation}\label{eq:lambdastar}
\lambda_*=\frac{(4\pi)^2}3\,\epsilon+O(\epsilon^2)\,,
\end{equation}
but for sufficiently small $\epsilon$ these corrections cannot destroy the fixed
point. The fixed point \eqref{eq:lambdastar} is the Wilson-Fisher fixed point.

The above result was obtained using renormalized perturbation theory for $d<4$
$\varphi^4$ theory. As I mentioned, for $d<4$ the unrenormalized perturbation
theory, i.e.\ the expansion in terms of the bare coupling constant $\lambda_0$, is
(essentially) finite. It is instructive to use this example to explain the
relation between renormalized and unrenormalized perturbation theories.

If we repeat our calculation of the $\beta$-function using unrenormalized
perturbation theory, i.e.\ with the expression \eqref{eq:G4bare} for the
four-point vertex, we would obtain from \eqref{eq:CSGamma21}
\begin{equation}\label{eq:beta0}
\beta_0(\lambda_0)=-\epsilon\,\lambda_0\,,
\end{equation}
with no fixed point visible. Higher order corrections do not change this result;
in fact, it is easy to see that in unrenormalized perturbation theory, without the
counterterm contributions, the formula \eqref{eq:beta0} is exact. If one uses
$\lambda_0$ as the coordinate on $U(\mathcal A_{*G})$, the RG function $\beta$
reduces to its linear part. The bare coupling constant $\lambda_0$ is the
``canonical coordinate'', whose existence is guaranteed by the Poincar\'e-Dulac
theorem.

Indeed, in the case $\epsilon\ll1$ the RG flow equation for small values of the
renormalized parameter $\lambda$ has the form
\begin{equation}\label{eq:logistic}
\frac{(4\pi)^2}3\,\frac{\dd}{\dd l}\,\lambda=\lambda\,(\lambda_*-\lambda)\,.
\end{equation}
It is not difficult to check that the change of variables
\begin{equation}\label{eq:lambda0map}
\lambda=\frac{\lambda_0}{1+\frac{\lambda_0}{\lambda_*}}
\end{equation}
linearizes this flow equation,
\begin{equation}\label{eq:lambda0flow}
\frac{\dd}{\dd l}\,\lambda_0=\epsilon\,\lambda_0\,.
\end{equation}
Note that the transformation \eqref{eq:lambda0map}, or equivalently
\[
\frac1\lambda=\frac1{\lambda_0}+\frac1{\lambda_*}\,,
\]
maps the interval $0<\lambda<\lambda_*$ to the whole positive half-axis of
$\lambda_0$,
\[
\lambda\in[0,\lambda_*)\quad\leftrightarrow\quad\lambda_0\in[0,\infty)\,,
\]
with the fixed point $\lambda_*$ mapped to the point $\lambda_0=\infty$; for this
reason it is invisible in \eqref{eq:lambda0flow}. Note also that if one expands
$\lambda$ in terms of $\lambda_0$ according to \eqref{eq:lambda0map}, for small
$\epsilon$ the coefficients of this expansion become large,
\[
\lambda=\lambda_0-\frac{\lambda_0^2}{\lambda_*}+\cdots\,,
\]
which explains the large coefficients in \eqref{eq:G4bare}. For $\epsilon=0$
(i.e.\ $d=4$) the interval $[0,\lambda_*)$ shrinks to zero, the perturbation
becomes marginal, and the ``canonical coordinate'' $\lambda_0$ ceases to exist.

If the renormalized parameter $\lambda$ is used, for small $\epsilon$ the flow
between $\mathcal A_{*G}$ and $\mathcal A_{*WF}$, as well as the fixed point
$\mathcal A_{*WF}$ itself, remains at small $\lambda\sim\epsilon$, within the
domain where the perturbative results are reliable. Therefore for small
$\epsilon$ we can use the renormalized perturbation theory to calculate all
characteristics of the Wilson-Fisher fixed point. For instance, to the leading
order in $\epsilon$ the flow equation for the mass parameter $M^2$ has the same
form as at $d=4$,
\begin{equation}\label{eq:M2flow21}
\frac{\dd}{\dd l}\,M^2=\left(2-\frac\lambda{(4\pi)^2}+O(\epsilon^2)\right)\,M^2\,.
\end{equation}
Together with \eqref{eq:logistic} this equation yields the pattern of the flow
shown in figure~\ref{fig:wfflow}.

\subsection*{\texorpdfstring{$\epsilon$}{epsilon}-expansion}

The main assumption of $\epsilon$-expansion is that the general features of the
RG flow observed at small $\epsilon$ remain qualitatively the same at
$\epsilon=1$, i.e.\ at $d=3$. Under this assumption, the characteristics of the
Wilson-Fisher fixed point can be calculated perturbatively, as the power series in
$\epsilon$. Although at $d=3$ the expansion parameter $\epsilon$ is not small,
this approach yields remarkably accurate results for the critical exponents in
three dimensions.

To find the critical exponents one must calculate the dimensions of various
relevant fields at the Wilson-Fisher fixed point. Note that $\Phi_0=\varphi^2$ is
the only symmetric relevant field at $\mathcal A_{*WF}$. The dimensions of the
fields $\Phi_0$ and $\varphi$ at this point can be calculated as
\[
D_0=D_{\varphi^2}=d-2+\gamma_{\varphi^2}(\lambda_*)\,,\qquad
D_\varphi=\frac{d-2}2+\gamma(\lambda_*)\,.
\]
In the leading order in $\epsilon$ we can just borrow the functions $\gamma$ from
the $d=4$ calculation,
\[
\gamma_{\varphi^2}(\lambda)=\frac\lambda{(4\pi)^2}+O(\lambda^2)\,,\qquad
\gamma(\lambda)=\frac1{12}\,\frac{\lambda^2}{(4\pi)^4}+O(\lambda^3)\,,
\]
and use the value \eqref{eq:lambdastar} for $\lambda_*$. This yields for the
critical exponents
\[
\frac1\nu=d-D_0=2-\frac\epsilon3+O(\epsilon^2)\,,\qquad
\alpha=2-\frac d{d-D_0}=\frac\epsilon6+O(\epsilon^2)\,,
\]
and similar results for other exponents. For $\epsilon=1$ this gives $\nu=0.60$
and $\alpha=0.16$ for the correlation length and heat capacity exponents, to be
compared to $\nu=0.63$ and $\alpha=0.12$ found in experiment.\footnote{The
definition is $R_c\sim|T_c-T|^{-\nu}$, $C\sim|T_c-T|^{-\alpha}$.}

\subsection*{\texorpdfstring{$O(N)$ $\varphi^4$}{O(N) phi**4} theory}

Let us discuss somewhat more general $\varphi^4$ theory. Assume that the scalar
field has $N$ components $\bm{\varphi}=(\varphi_1,\dots,\varphi_N)$, and the
$\varphi^4$ interaction respects the $O(N)$ transformations of the vector
$\bm{\varphi}$. We have
\[
\mathcal A=\int \dd^dx\,\left[\frac12\,(\partial\bm{\varphi})^2+\frac{M^2}2\,\bm{\varphi}^2+\frac\lambda{4!}\,(\bm{\varphi}^2)^2\right]
\]
Although, as before, we use the notation $M^2/2$ for the coefficient in front of
the term $\bm{\varphi}^2$, this coefficient is allowed to take negative values.

At small $d=4-\epsilon$ with small $\epsilon$ the pattern of the RG flow of the
$O(N)$ $\varphi^4$ theory in the vicinity of its Gaussian fixed point is
essentially the same as in the one-component case we discussed
previously\footnote{Explicit calculation (Problem~\ref{hw:8}) yields, in the
normalization $\frac\lambda{4!}(\bm{\varphi}^2)^2$ used here,
$\beta(\lambda)=-\epsilon\lambda+\frac{N+8}{48\pi^2}\,\lambda^2+\cdots$,
$\gamma_{\varphi^2}(\lambda)=\frac{N+2}{48\pi^2}\,\lambda+\cdots$.}
(figure~\ref{fig:onflow}).
\begin{figure}[ht]
\centering\figfile{onflow}
\caption{RG flow of the $O(N)$ $\varphi^4$ theory near $d=4$. The axis $M^2=0$ is
the critical surface; above it the flow goes up, below it down.}
\label{fig:onflow}
\end{figure}

There is the $O(N)$ version of the Wilson-Fisher fixed point
$\mathcal A_{*,WF}\in\Sigma_c$ at $M^2=0$, $\lambda=\lambda_*\sim\epsilon$. In the
space of $O(N)$ invariant interactions, it has one-dimensional unstable manifold
$U(\mathcal A_{*,WF})$.

In figure~\ref{fig:onflow} the axis $M^2=0$ is the critical surface $\Sigma_c$.
Above it, the flow generally goes upward, while everything below it flows
downward. One can try to figure out the ultimate fate of the flow at positive and
negative $M^2$ on physical grounds. The domain $M^2>0$ corresponds to the
high-temperature phase of the ferromagnetic system. The pattern of the flow shows
that $M^2$ increases with $l$. At large positive $M^2$ the term
$\frac{M^2}2\bm{\varphi}^2$ suppresses fluctuations of $\bm{\varphi}$. Therefore
at $M^2>0$ one expects to have attractive non-critical fixed point
$P_{\text{high}}$, corresponding to the probability distribution
\[
P_{\text{high}}[\bm{\varphi}]\sim\prod_x\delta\bigl(\bm{\varphi}(x)\bigr)\,.
\]
The domain of negative $M^2$ is associated with the low-temperature phase. Here
under the RG flow $M^2$ remains negative while increasing in the absolute value. In
this domain the potential term
$\frac{M^2}2\bm{\varphi}^2+\frac\lambda{4!}(\bm{\varphi}^2)^2$, as the
function of $\bm{\varphi}^2$, has the minimum away from zero
(figure~\ref{fig:mexican}).
\begin{figure}[ht]
\centering\figfile{mexican}
\caption{The potential $V$ as a function of $|\bm{\varphi}|$ for $M^2<0$.}
\label{fig:mexican}
\end{figure}

In terms of $\bm{\varphi}$ we have manifold of degenerate minima, the sphere
$S^{N-1}$, of the radius
\[
|\bm{\varphi}|=\rho:=\sqrt{\frac{6\,|M^2|}\lambda}\,.
\]
Points of the sphere indicate possible directions of $\bm{\varphi}$ at the minimum.
It is easy to see that if $|M^2|$ is large the minimum in figure~\ref{fig:mexican}
becomes deep and narrow. Fluctuations of the radial part, $|\bm{\varphi}|$, are
then suppressed. On the other hand there is no suppression for the fluctuations of
the angular part. Writing
\[
\bm{\varphi}=\rho\,\mathbf n\,,\qquad|\mathbf n|=1\,.
\]
If we completely ignore the radial fluctuations, i.e.\ set $|\bm{\varphi}|=\rho$
everywhere, the action takes simple form
\[
\mathcal A=\int \dd^dx\,\left[\frac1{2g}\,(\partial\mathbf n)^2\right]
\qquad\text{with}\qquad
g=\frac1{\rho^2}=\frac\lambda{6\,|M^2|}\,.
\]
As $|M^2|$ increases with the RG flow, $g$ becomes small.\footnote{This simple
classical analysis is valid only in the limit $\lambda\to0$. More generally,
fluctuations of $|\bm{\varphi}|$ can be taken into account, along with the
fluctuations of the angular part, by renormalized perturbation theory. Some
details about this theory, the so called linear sigma model, can be found in
Chapter 11 of PS.} Due to the non-linear constraint $\mathbf n^2=1$ the above
action leads to a theory more complicated than a free massless field. It is named
the non-linear sigma model.

Let us analyze physical content of the nonlinear sigma-model. Classically, the
most favorable field configurations correspond to minima of the action, which in
this case are constant ($x$-independent) fields
\[
\mathbf n(x)=\mathbf n_0\,.
\]
By the $O(N)$ symmetry we can choose
\[
\mathbf n_0=(0,0,\dots,0,1)
\]
with no loss of generality. To take into account fluctuations around this minimum
one can set
\[
\mathbf n(x)=\bigl(\pi^1(x),\pi^2(x),\dots,\pi^{N-1}(x),\sigma(x)\bigr)\equiv\bigl(\pi^k(x),\sigma(x)\bigr)\qquad(k=1,2,\dots,N-1)
\]
with
\[
\sigma(x)=\sqrt{1-(\pi^k(x))^2}
\]
The action takes the form
\[
\mathcal A=\int \dd^dx\,\frac1{2g}\left[(\partial_\mu\pi^k)^2+\frac{(\pi^k\partial_\mu\pi^k)^2}{1-(\pi^k)^2}\right]
\]
Ignoring first the quartic and higher terms in this action, we observe $N-1$
massless scalar fields $\pi^k(x)$. These are the \emph{Goldstone bosons},
associated with the spontaneous breakdown of the $O(N)$ symmetry, down to
$O(N-1)$ (generalities about the spontaneous symmetry breaking are in Chapter 11
of PS).

Rescaling
\[
\pi^k=\sqrt g\,\phi^k
\]
one can write the above action as
\[
\mathcal A=\int \dd^dx\,\left[\frac12\,(\partial\phi^k)^2+\sum_{n=1}^\infty g^n\,O_n(x)\right]\,,
\qquad\text{where}\qquad
O_n=\frac12\,(\phi^k\partial_\mu\phi^k)^2\,(\phi^k\phi^k)^{n-1}\,.
\]
If $d>2$, it is correct to think of it as $N-1$-component of massless free field
(the first term), with some composite fields added. The dimensions of the fields
$O_n$ are
\[
O_n:\qquad D_n=2+(n+1)(d-2)=d+n\,(d-2)\,.
\]
At $d>2$ all these fields are irrelevant.

By these arguments, at $d>2$ the long-distance physics in the low-temperature
domain $M^2<0$ is controlled by the action
\[
\mathcal A_{*,GS}=\int\frac12\,(\partial\phi^k)^2\dd^dx
\]
which is essentially a Gaussian fixed point of $N-1$ massless bosons. I will call
it the ``Goldstone point'', to distinguish it from the original Gaussian fixed
point $\mathcal A_{*,G}$ (which has $N$ massless scalar fields).

Returning to the pattern of the RG flow on $U(\mathcal A_{*,G})$, we now anticipate
the global structure shown in figure~\ref{fig:onglobal}.
\begin{figure}[ht]
\centering\figfile{onglobal}
\caption{Anticipated global RG flow on $U(\mathcal A_{*,G})$ for $N>1$: the flow
above the critical surface ends at $P_{\text{high}}$, and the flow below it at the
critical Goldstone point $\mathcal A_{*,GS}$.}
\label{fig:onglobal}
\end{figure}
For $N>1$ there is the \emph{critical} fixed point $\mathcal A_{*,GS}$, which
attracts the trajectories in the low-temperature domain $M^2<0$.

At $d=2$ we have $D_n=2=d$, that is all the above fields $O_n$ become marginal.
Obviously, something interesting is happening in this case. Let us analyze it
separately.

\begin{hwproblem}{8}\label{hw:8}
Consider the same $O(N)$ $(\bm{\varphi})^4$ theory as in the HW
Problem~\ref{hw:6}, but this time in $d=4-\epsilon$ space, with $\epsilon\ll1$.

Find the Wilson-Fisher fixed point $\lambda_*$ to the leading order in $\epsilon$,
and evaluate the anomalous dimensions of the fields $(\bm{\varphi})^2$ and
$\bm{\varphi}$ to the order $\sim\epsilon$, at the WF fixed point. Find the critical
exponents $\nu$ in $R_c\sim|T_c-T|^{-\nu}$, and $\alpha$ in
$C\sim|T_c-T|^{-\alpha}$.
\end{hwproblem}
